\subsection{绝对半单代数}

定义 2. —— 设 K 是一个交换域。我们称 K-代数 A 是绝对半单的，如果对 K 的每个扩张 L，环 \(\mathrm{A}_{(\mathrm{L})}\) 都是半单的。

绝对半单代数是半单的。K-代数 A 绝对半单当且仅当 A-模 \(A_{s}\) 绝对半单。于是，由 VIII, 第 229 页的命题 1，我们得到如下结果：若

L 是 K 的一个扩张，则 L-代数 \(A_{(L)}\) 绝对半单当且仅当 K-代数 A 绝对半单。

定理 1. —— 设 K 是一个交换域，A 是一个 K-代数。下列性质等价：

\begin{enumerate}
\def\labelenumi{(\roman{enumi})}
\item
  K-代数 A 是绝对半单的。
\item
  代数 \(\mathrm{A}\) 在 \(\mathrm{K}\) 上的次数有限，且存在 \(\mathrm{P}\)（\(\mathrm{K}\) 的一个为完满域的扩张），使得 \(\mathrm{P}\) -代数 \(\mathrm{A}_{(\mathrm{P})}\) 是半单的。
\item
  K-代数 A 是半单的且在 K 上的次数有限，并且其中心是一个平展 K-代数。
\item
  存在一个有限族 \((n_{i},\mathrm{D}_{i})_{i\in\mathrm{I}}\)，其中 \(n_{i}\) 是严格正整数，而 \(D_{i}\) 是 K 上的一个为体的有限次数代数，使得 \(Z_{i}\)（即 \(D_{i}\) 的中心）对每个 \(i\in I\) 都是 K 的一个可分扩张，且 A 同构于矩阵环 \(\mathbf{M}_{n_{i}}(\mathrm{D}_{i})\) 之积。
\item
  存在 \(\mathrm{L}\)（\(\mathrm{K}\) 的一个扩张）与一个有限整数族 \((n_i)_{i\in \mathrm{I}}\)，使得 L-代数 \(\mathrm{A}_{(\mathrm{L})}\) 同构于代数 \(\prod_{i\in \mathrm{I}}\mathbf{M}_{n_i}(\mathrm{L})\)。
\item
  存在 K 的一个有限次数伽罗瓦扩张 L 与一个有限整数族 \((n_{i})\)，使得 \(\mathrm{A}_{(\mathrm{L})}\) 同构于矩阵代数 \(\mathbf{M}_{n_{i}}(\mathrm{L})\) 之积。
\end{enumerate}

我们先证明诸蕴含关系 (v) \(\Rightarrow\) (iv) \(\Rightarrow\) (iii) \(\Rightarrow\) (ii) \(\Rightarrow\) (i)。用 Z 表示 A 的中心。

若性质 (v) 成立，则 L-代数 \(\mathrm{A}_{(\mathrm{L})}\) 是半单的且在 L 上的次数有限，其中心（同构于 \(\mathrm{Z}_{(\mathrm{L})}\)，见 III, 第 41 页的推论）同构于 \(L^{r}\)，其中 \(r \geqslant 0\) 是某个整数。由 VIII, 第 222 页的推论 2, a)，代数 A 是半单的且在 K 上的次数有限。因此它同构于环的一个有限乘积 \(\prod_{i \in I} M_{n_{i}}(D_{i})\)，其中 \(n_{i} \geqslant 1\) 对每个 \(i \in I\) 成立，而体 \(D_{i}\) 是 K 上的一个有限次数代数。\(Z_{i}\)（即 \(D_{i}\) 的中心）是一个交换域且是 K 的一个扩张，而 Z 同构于 \(\prod_{i \in I} Z_{i}\)。于是，\(\mathrm{Z}_{(\mathrm{L})}\) 一方面同构于 \(\prod_{i \in I} (\mathrm{Z}_{i})_{(\mathrm{L})}\)，另一方面同构于 \(L^{r}\)。换言之，代数 \(\prod_{i \in I} Z_{i}\) 在域 K 上是平展的（V, §6, No.~3, 第 28 页，定义 1），且诸扩张 \(Z_{i}\) 在 K 上都是可分的（V, §6, No.~4, 第 31 页，推论，及 V, §7, No.~1, 第 42 页）。故 (v) 蕴含 (iv)。

若性质 (iv) 成立，则显然 A 是一个半单代数且在 K 上的次数有限。其中心 Z 同构于 K 的有限次数可分扩张的乘积 \(\prod_{i\in I}Z_{i}\)；因而它是一个平展代数（同前所引）。故 (iv) 蕴含 (iii)。

蕴含关系 (iii) \(\Rightarrow\) (ii) \(\Rightarrow\) (i) 由应用到 A-模 \(\mathrm{A}_s\) 的 VIII, 第 230 页的命题 2 得出。

引理 2. —— 设 L 是一个代数闭域，D 是一个在其中心中包含 L 的体。若 D 异于 L，则存在 L 的一个扩张 \(L'\)，使得环 \(D \otimes_{L} L'\) 不是左阿廷的。

设 \(x\) 是 \(\mathrm{D} - \mathrm{L}\) 的一个元素；由于 \(\mathrm{L}\) 是代数闭的，扩张 \(\mathrm{L}' = \mathrm{L}(x)\)（\(\mathrm{L}\) 的扩张）不是代数的，且 \(x\) 在 \(\mathrm{L}\) 上是超越的。于是环 \(\mathrm{B} = \mathrm{L}' \otimes_{\mathrm{L}} \mathrm{L}'\) 是一个整环，这由 \(\mathrm{V}\), §17, No.~4, 第 141 页的命题 5 得出。

B 的元素 \(y = x \otimes 1 - 1 \otimes x\) 不是零，但若 \(\varphi\) 是从 B 到 \(L'\) 的把 \(\xi \otimes \eta\) 映为 \(\xi \eta\) 的同态，则我们有 \(\varphi(y) = 0\)，故 y 在 B 中不可逆。我们把环 \(C = D \otimes_{L} L'\) 看作它的子环 B 上的右模；它是一个自由模，因为 D 是它的子域 \(L'\) 上的一个右向量空间。由于 y 是整环 B 的非零且不可逆的元素，在 C 中用 y 作右乘是一个单但非双射的映射 \(R_y\)。而 \(R_y\) 是左 C-模 \(C_s\) 的一个自同态；因此（VIII, 第 28 页，推论 1），环 C 不是左阿廷的。

我们现在来证明 (i) 蕴含 (v)。这是下述引理的一个推论。

引理 3. —— 设 A 是一个绝对半单 K-代数，L 是 K 的一个代数闭扩张。则代数 \(\mathrm{A}_{(\mathrm{L})}\) 同构于 L 上有限多个矩阵代数的一个乘积。

L-代数 \(\mathrm{A}_{(\mathrm{L})}\) 是半单的；因此它同构于有限多个形如 \(\mathbf{M}_{n_{i}}(\mathrm{D}_{i})\) 的代数的乘积，其中 \(D_{i}\) 是一个在其中心中包含 L 的体，而 \(n_{i}\) 是整数 \(\geqslant 1\)（VIII, 第 137 页，注 1）。

设 \(L'\) 是 L 的一个扩张。由于 K-代数 A 是绝对半单的，环 \(A_{(L')}\) 是半单的，因而是左阿廷的。现在，环 \(A_{(L')}\) 同构于 \(L' \otimes_L A_{(L)}\)，从而同构于 \(\prod_{i \in I} M_{n_i}(L' \otimes_L D_i)\)；由 VIII, 第 7 页的命题 5，诸环 \(M_{n_i}(L' \otimes_L D_i)\) 因而都是左阿廷的。

设 \(n \geqslant 1\) 是一个整数，B 是一个环。设 \((\mathfrak{b}_r)_{r \geqslant 0}\) 是 B 的左理想的一个递降序列；用 \(\mathfrak{c}_r\) 表示 \(n\) 阶且元素在 \(\mathfrak{b}_r\) 中的方阵所成之集。则 \((\mathfrak{c}_r)_{r \geqslant 0}\) 是 \(\mathbf{M}_n(\mathrm{B})\) 的左理想的一个递降序列。特别地，若环 \(\mathbf{M}_n(\mathrm{B})\) 是左阿廷的，则 B 也是。

由上所述，对每个 \(i \in I\) 与 L 的每个扩张 \(L'\)，环 \(D_{i} \otimes_{L} L'\) 都是左阿廷的。由引理 2，我们有 \(D_{i} = L\)（对每个 \(i \in I\)），由此即得引理 3。

我们将用下面的引理来证明蕴含关系 (i)⇒(vi)。

引理 4. —— 设 A 与 B 是域 K 上的具有有限生成系的代数，\(K'\) 是 K 的一个扩张。若 \(K'\) -代数 \(\mathrm{A}_{(\mathrm{K}')}\) 与 \(\mathrm{B}_{(\mathrm{K}')}\) 同构，则存在 \(K'\) 的一个在 K 上有限生成的子扩张 L，使得 L-代数 \(\mathrm{A}_{(\mathrm{L})}\) 与 \(\mathrm{B}_{(\mathrm{L})}\) 同构。

设 \((e_{i})_{i\in\mathrm{I}}\) 与 \((f_{j})_{j\in\mathrm{J}}\) 分别是代数 A 与 B 的有限生成系。设 u 是从 \(\mathrm{A}_{(\mathrm{K}^{\prime})}\) 到 \(\mathrm{B}_{(\mathrm{K}^{\prime})}\) 的一个同构，v 是其逆同构；存在 \(K^{\prime}\) 的一个在 K 上有限生成的子扩张 L，使得 \(u(1\otimes e_{i})\in\mathrm{B}_{(\mathrm{L})}\)（对每个 \(i\in\mathrm{I}\)）且 \(v(1\otimes f_{j})\in\mathrm{A}_{(\mathrm{L})}\)（对每个 \(j\in\mathrm{J}\)）。于是，u 把 \(\mathrm{A}_{(\mathrm{L})}\) 映到 \(\mathrm{B}_{(\mathrm{L})}\) 中，而 v 把 \(\mathrm{B}_{(\mathrm{L})}\) 映到 \(\mathrm{A}_{(\mathrm{L})}\) 中。诱导映射 \(u^{\prime}:\mathrm{A}_{(\mathrm{L})}\to\mathrm{B}_{(\mathrm{L})}\) 与 \(v^{\prime}:\mathrm{B}_{(\mathrm{L})}\to\mathrm{A}_{(\mathrm{L})}\) 是环同态；它们是互逆的双射。

我们来完成蕴含关系 \((\mathrm{i})\Rightarrow(\mathrm{vi})\) 的证明。设 \(K'\) 是 K 的一个可分闭包（V, §7, No.~8, 第 45 页）。则 \(A_{(K')}\) 是 \(K'\) 上的一个绝对半单代数。由蕴含关系 \((\mathrm{i})\Rightarrow(\mathrm{iv})\)，\(K'\) -代数 \(A_{(K')}\) 同构于一个乘积 \(\mathbf{M}_{n_{1}}(\mathrm{D}_{1})\times\cdots\times\mathbf{M}_{n_{r}}(\mathrm{D}_{r})\)，其中 \(D_{i}\) 是 \(K'\) 上的一个为体的有限次数代数，其中心 \(Z_{i}\) 是 \(K'\) 的一个可分扩张。由于 \(K'\) 是可分闭的，我们有 \(Z_{i}=K'\)。由 VIII, 第 231 页的命题 3，体 \(D_{i}\) 是交换的。因此我们有 \(D_{i}=K'\)。用 B 表示 K-代数 \(\mathbf{M}_{n_{1}}(\mathrm{K})\times\cdots\times\mathbf{M}_{n_{r}}(\mathrm{K})\)。由上所述，\(K'\) -代数 \(A_{(K')}\) 与 \(B_{(K')}\) 同构。\(K'\) 的每个在 K 上有限生成的子扩张都是可分的且在 K 上的次数有限，因此，每一个这样的子扩张都包含于某个子扩张 L 之中，其中 L 是 \(K'\) 的在 K 上伽罗瓦且次数有限的子扩张（V, §10, No.~1, 第 57 页，命题 2）。于是该蕴含关系由引理 4 得出。

蕴含关系 (vi)⇒(v) 是直接的。

推论 1. —— 设 K 是一个交换域，\(A_{1}\) 与 \(A_{2}\) 是两个 K-代数。假设 \(A_{1}\) 是绝对半单的。

\begin{enumerate}
\def\labelenumi{\alph{enumi})}
\item
  若 \(\mathrm{A}_2\) 是半单的，则 \(\mathrm{A}_1\otimes_{\mathrm{K}}\mathrm{A}_2\) 亦然
\item
  若 \(\mathrm{A}_2\) 是绝对半单的，则 \(\mathrm{A}_1\otimes_{\mathrm{K}}\mathrm{A}_2\) 亦然。
\end{enumerate}

把 \(A_{1}\) 的中心记作 \(Z_{1}\)，把 \(A_{2}\) 的中心记作 \(Z_{2}\)。由 III, §4, No.~4, 第 468 页的推论，\(A_{1} \otimes_{K} A_{2}\) 的中心 Z 等于 \(Z_{1} \otimes_{K} Z_{2}\)。假设 \(A_{2}\) 是半单的；则 \(Z_{2}\) 是一个约化代数（VIII, 第 137 页，命题 2 与 3）。由定理 1，\(Z_{1}\) 是平展的，因而是可分的

K-代数。由 V, §15, No.~2, 第 120 页的命题 5，环 \(Z = Z_{1} \otimes_{K} Z_{2}\) 是约化的；由于 \(A_{1}\) 在 K 上的次数有限（定理 1），由 VIII, 第 221 页的命题 7 得出，环 \(A_{1} \otimes_{K} A_{2}\) 是半单的。

现在假设 \(A_{2}\) 是绝对半单的。设 L 是 K 的一个扩张。则代数 \(\mathrm{A}_{1(\mathrm{L})}\) 是绝对半单的，而代数 \(\mathrm{A}_{2(\mathrm{L})}\) 是半单的。因此由 a)，代数 \(\mathrm{A}_{1} \otimes_{\mathrm{K}} \mathrm{A}_{2(\mathrm{L})}\) 是半单的。

推论 2. —— 设 K 是一个可分闭域，A 是一个绝对半单 K-代数。则存在整数 \(r \geqslant 0\) 与严格正整数 \(n_{1}, \ldots, n_{r}\)，使得代数 A 同构于代数 \(\prod_{i=1}^{r} \mathbf{M}_{n_{i}}(\mathrm{K})\)。

由定理 1，A 同构于形如 \(\prod_{i=1}^{r} \mathbf{M}_{n_i}(\mathrm{D}_i)\) 的一个代数，其中 \(r \geqslant 0\) 是整数，\(n_1, \ldots, n_r\) 是整数，而 \(\mathrm{D}_1, \ldots, \mathrm{D}_r\) 是 K 上的为体的有限次数代数，它们的中心是 K 的可分扩张，因而等于 K。由 VIII, 第 231 页的命题 3，我们有 \(\mathrm{D}_i = \mathrm{K}\)（对 \(i \in [1, r]\)）。

例. —— 一个交换 K-代数是绝对半单的当且仅当它是平展的：这由定义（V, §6, No.~3, 第 28 页，定义 1）与定理 1 的性质 (i) 与 (v) 的等价性得出。

